\section{考察}
\subsection{従来手法との差別化点}
本研究の差別化点は，以下の3点に整理できる．

第一に，YOLO\cite{ref_YOLO}やDETR\cite{ref_DETR}のような閉域的な物体検出モデルに対しては，学習時に定めた固定のクラス集合を超えたカテゴリを予測できないという制約を解消している．
本手法では，DETRの予測カテゴリを用いず，バウンディングボックス出力のみを物体領域の検出に用いることで，検出済みの領域に対し後段のVLMが任意の候補語彙でカテゴリをアノテーションできるようにしている．
4.6節で確認したとおり，この設計により，COCOの80クラスに含まれない語（penguin，helicopter，kangaroo，camel）についても，追加学習なしにアノテーションが可能である．
\add{ただし，4.8節で示したとおり，COCOの80クラスの範囲内では，そのクラスで教師あり学習したYOLOの方がカテゴリのアノテーション精度は大幅に高い．本手法の利点は学習時の語彙を超えた語や粒度の異なる語でアノテーションできる点にあり，閉域の語彙で十分な用途ではYOLOのような閉域モデルが適している．}

第二に，OWL-ViT\cite{ref_OWLViT}やGrounding DINO\cite{ref_GDINO}のようなオープンボキャブラリ物体検出モデルに対しては，候補カテゴリ集合を変更するたびに検出処理全体を再実行する必要があるという制約を解消している．
4.3節で確認したとおり，位置検出精度はDETRとこれらのモデルの間で大きな差がないため，本手法は同等の位置検出精度を維持しつつ，物体領域の検出を1回のみに抑え，後段のカテゴリのアノテーションのみを候補語彙集合の変更に応じて繰り返す設計を実現している．

第三に，ImageNetのクラス階層を利用する従来の階層的画像分類の多くは，あらかじめ定めた固定の階層構造に対して学習を行うものであり，推論時に階層の深さや分岐を動的に選択することは想定していない．
本手法は，WordNetの上位語・下位語関係を推論時にたどることで，学習を一切行わずに粒度可変のアノテーションを実現している点で異なる．

\subsection{クロップの前処理と検出スコアの閾値の影響}
4.4節の結果は，同一のDETRとVLMを用いても，クロップの前処理と検出スコアの閾値によってTop-1正解率が0.481から0.604まで変化することを示している．
特に前処理の影響は，padでもcenterでも物体の外側の画像内容は含まないにもかかわらず，VLMの標準前処理である中央切り出しによって細長い物体の一部が失われることに起因する．
高さが幅の2倍以上のクロップでは，$ \theta = 0.9 $においてTop-1正解率がcenterの0.441からpadの0.545へ向上した．
クロップを入力とするVLMのゼロショット分類では，前処理を明示することが精度の比較に不可欠である．
一方，人物クラスでは前処理によらず正解率が0.38程度にとどまった．
人物のクロップには，人物とともにバット，ネクタイ，ラケット等の物体が含まれることが多く，``a photo of a person''よりもそれらの物体のプロンプトとの類似度が高くなる．
これは，候補語彙の中で包含関係にある物体同士をどのように扱うかという，クロップに基づくアノテーションに固有の課題である．

\subsection{WordNet階層をそのまま用いることの限界}
4.5節，4.7節の結果は，WordNetの生の階層構造をそのまま候補語彙として用いることに起因する2種類の限界を示している．

第一に，深さ0から深さ1への確信度の大幅な低下（0.642から0.442）に見られるように，候補語彙数の増加および視覚的に類似した細かいカテゴリ間の識別困難性により，最初の細分化の段階でアノテーションの信頼性が低下する．
さらに，\texttt{person.n.01}の下位語には``suspect''（20件），``thin person''（24件），``visually impaired person''（9件）等，外見だけからは判断できない社会的・属性的な語が多く含まれ，これらが実際に選択されていた．

第二に，giraffe → okapiのように，WordNetの階層構造には，視覚的なカテゴリ分類の候補語彙として不適切な分岐が含まれている．
特に下位語が1語しかない分岐では，(\ref{eq_clip})のソフトマックスにより確信度が画像によらず1となり，候補選択が実質的に機能しない．
このような段階は細分化の段階全体の約18\%を占め，深さ2以降の平均確信度を見かけ上押し上げていた．
単一候補の段階を除くと，深さ1から深さ3の平均確信度は0.378，0.571，0.481であり，深さ0と比べて一貫して低い．
これらの限界は，WordNetの階層構造が視覚的な類似性ではなく言語学的な意味関係に基づいて構築されていることに起因すると考えられる．

\subsection{物体領域の検出のCOCO学習分布への依存}
4.6節の灯台の事例は，カテゴリに依存しない検出器として利用する場合であっても，DETRの物体領域の検出そのものがCOCOの学習データの分布から完全には独立ではないことを示している．
すなわち，本手法における「物体領域の検出とカテゴリのアノテーションの分離」は，カテゴリの予測に関する制約を解消するものであるが，物体領域の候補生成そのものに関する制約は解消しない．
学習データに存在する物体の形状・構図との類似性への依存は残る．
また，ヘリコプターの吊り荷が``backpack''として誤ってアノテーションされた事例は，候補語彙集合に正解語が存在しない領域に対して，VLMが候補語彙中で最も類似度の高い無関係な語を選択してしまう問題を示している．

\subsection{今後の課題}
今後の課題として，第一に，WordNet階層のフィルタリングによる粒度可変アノテーションの実用性向上が挙げられる．
具体的には，下位語が1語しかない分岐の打ち切りや，絶滅種・希少種・外見から判断できない語の除外である．
第二に，物体領域の検出モデルの比較を本稿の4画像から500枚規模へ拡大し，統計的な比較を行うことが挙げられる．
第三に，物体領域の検出における学習分布への依存を緩和するため，Segment Anything Model（SAM）等の汎用的なセグメンテーションモデルとDETRを組み合わせるなど，物体領域の検出モデル自体の改善が挙げられる．
第四に，候補語彙に正解語が存在しない領域を棄却する仕組みとして，確信度に対する閾値処理の高度化や，候補語彙に「その他」に相当する語を加える方法の導入が挙げられる．
第五に，人物クロップにおける持ち物・装着物への誤反応を抑えるため，プロンプトの設計や，検出済みの他の領域との包含関係を考慮したアノテーションが挙げられる．
\add{第六に，4.8節の結果を踏まえ，COCOの80クラスで高精度なYOLOの予測を深さ0のアノテーションとして用い，深さ1以降の細分化と80クラスに含まれない語のアノテーションのみをVLMで行う組合せが挙げられる．}
